\newpage
\section{Conclusion and Future Work}
\label{sec:conclusion}

\noindent This study introduced and evaluated an A* search-based framework for training Multilayer Perceptrons, leveraging a non-admissible heuristic function to navigate a quantized parameter space. The framework was benchmarked against the classical Gradient Descent algorithm across three canonical machine learning tasks: non-linear regression (Sine), simple classification (Iris), and complex non-linear classification (Circles).

\subsection{Summary of A* Framework Performance}
\label{sec:astar_summary}

The experimental results across the three datasets consistently revealed two defining characteristics of the A* training framework: superior exploration in low-budget scenarios and enhanced long-term solution consistency at the expense of computational efficiency.

\subsubsection*{1. A* as a Superior Low-Budget Explorer (Sine Dataset)}

On the non-linear regression task (Sine), A* demonstrated a clear advantage under a limited budget ($I_{\text{max}}=1,000$).
\begin{itemize}
    \item \textbf{Central Tendency:} A* achieved a significantly lower Average Loss ($0.2852$) and Median Loss ($0.3319$) compared to GD Average Loss ($0.3852$) and Median Loss ($0.3843$).
    \item \textbf{Global Search:} A* showcased superior Best-Case Potential, finding solutions with a Minimum Loss ($L_{\min} = 0.0575$) that were four times better than GD's best. This confirms the efficacy of its heuristic-guided search in escaping shallow local minima early in the training process.
\end{itemize}
This evidence suggests that in situations where computational time is not a constraint, the A* framework can yield a higher-quality result than GD.

\subsubsection*{2. A* as a Consistent and Robust Trainer (Sine and Circles Datasets)}

While GD eventually surpassed A* in average performance in high-budget scenarios (Sine at $I_{\text{max}}=5,000$ and Circles at $I_{\text{max}}=2,500$), A* consistently offered a critical advantage: superior consistency (low variance).
\begin{itemize}
    \item \textbf{Consistency Reversal (Sine):} As the budget increased from $1,000$ to $5,000$ iterations, GD's consistency dramatically worsened ($\sigma$ increased), while A*'s consistency improved ($\sigma$ decreased). At $I_{\text{max}}=5,000$, A* was demonstrably more consistent ($\sigma = 0.1243$) than GD ($\sigma = 0.1603$).
    \item \textbf{Non-Linear Robustness (Circles):} Even in the challenging non-linear Circle dataset at $I_{\text{max}}=2,500$, where GD achieved a better median, A* maintained a lower Standard Deviation ($\sigma = 0.2129$) and Variance ($\sigma^2 = 0.0453$), leading to a better overall Average Loss ($\bar{L}=0.1703$).
\end{itemize}

\noindent This consistent behavior suggests that the A* framework provides a more reliable guarantee of solution quality due to its lower relative variance compared to GD. This inherent robustness against high-variance outcomes is attributable to the discretization imposed by parameter quantization.

\noindent \textbf{The Role of Quantization in Variance Reduction:} The A* method operates within a finite, reduced-cardinality search space defined by the quantization factor ($Q$) and the parameter range. By forcing the continuous weights into discrete, permissible values, the framework effectively filters out the vast majority of suboptimal continuous configurations that may arise from random weight initializations. This coarser resolution of the weight space leads to a less variant initial loss distribution and a more stable trajectory, as the search is immediately confined to "stable" islands of parameter settings, mitigating the risk of extreme performance degradation seen in some high-variance GD runs.

\noindent \textbf{The Role of the Heuristic in Exploration:} Furthermore, the heuristic contributes to controlled exploration. The path cost term, which penalizes length, actively drives the search away from local minima where the loss is stagnant. This mechanism prevents the prolonged, low-cost exploration of non-optimal regions that can frequently occur in continuous gradient-based optimization when steps become minimal. Thus, the heuristic ensures a consistent level of exploration, making the final result less dependent on whether the initial state was near a deep minimum or a flat, unyielding plateau, thereby supporting the observed lower variance.

\subsubsection*{3. Limitations: Computational Cost}

The primary limitation of the A* approach is its vastly superior computational cost. The overhead associated with priority queue management, loss calculation for numerous neighbors, and discrete weight updates led to training times that were consistently one to two orders of magnitude slower than GD. Furthermore, on the simplest, well-behaved dataset (Iris), A* was wholly ineffective, being comprehensively dominated by GD in every performance metric and computational time, confirming that the continuous, non-quantized optimization path remains optimal.

\subsection{Summary of Comparative Findings}
\label{sec:comparative_findings}

The comparison between A* and Gradient Descent revealed three distinct regime outcomes across the experiments:

\begin{itemize}
    \item \textbf{Low Budget / Complex Landscape (Sine, $I_{\text{max}}=1k$):} A* is the performance leader, achieving significantly lower average loss and minimum loss, confirming its superior ability for rapid exploration and discovery of high-quality solutions.
    \item \textbf{High Budget / Complex Landscape (Sine, $I_{\text{max}}=5k$; Circles, $I_{\text{max}}=2.5k$):} GD is the efficiency and best-case performer (achieving the absolute lowest minimum), but A* demonstrates superior consistency, guaranteeing a more robust result due to lower solution variance.
    \item \textbf{Low Budget / Simple Landscape (Iris, $I_{\text{max}}=1k$):} For the short budget, Gradient Descent is the dominant method, achieving superior performance across all metrics (speed, average loss, minimum loss, and consistency) compared to A*. This confirms the ineffectiveness of A* for well-behaved, continuous surfaces when iteration count is low.
    \item \textbf{High Budget / Simple Landscape (Iris, $I_{\text{max}}=2.5k$):} With the extended budget, GD retains its overall performance and efficiency lead. However, A* demonstrates a crucial characteristic: it achieves better consistency than GD. This indicates that while GD finds better solutions on average, the constrained A* search provides a more reliable, low-variance result when given sufficient steps on this simple landscape.
\end{itemize}

\subsection{Future Work}
\label{sec:future_work}

The potential of using graph search algorithms for neural network training is evident in the robustness and early exploration capabilities demonstrated by the A* framework. Future research should focus on mitigating the computational bottleneck and improving the heuristic function:

\begin{enumerate}
    \item \textbf{Pruning and Heuristic Improvement:} Developing a faster, admissible heuristic and implementing aggressive search space pruning techniques (e.g., beam search, iterative deepening) is crucial to reduce the neighbor evaluation overhead.
    \item \textbf{Dynamic Quantization:} Implementing a system where the quantization factor $Q$ or the parameter range is dynamically adjusted during training could allow the framework to benefit from the coarse-grained exploration of A* initially, followed by the fine-grained precision of GD or a finer quantization later.
    \item \textbf{Hybrid Optimization:} The most promising direction is a Hybrid Optimizer that uses A* for the first $10\%-20\%$ of training (to locate an excellent region in the quantized space) and then switches to highly optimized Gradient Descent for the remaining fine-tuning.
    \item \textbf{Parallelization of Search Operations:} Investigating the possibility of parallelizing the neighbor generation and evaluation processes (e.g., using multi-threading or GPU acceleration). Since the evaluation of the loss for each potential successor node is computationally independent, this approach could significantly reduce the high latency associated with the A* framework's primary computational bottleneck.
\end{enumerate}